% Part: computability
% Chapter: recursive-functions
% Section: partial-functions

\documentclass[../../../include/open-logic-section]{subfiles}

\begin{document}

\olfileid{cmp}{rec}{par}
\olsection{جزوي بازګشتي تابعې}


د بازګشتي تابعو د تعريف د انګېزې دپاره پام وکړئ چې زموږ هغه ثبوت، چې
محاسبه کېدونکې خو نابنسټيزې بازګشتي تابعې شته، په حقيقت کښې تر دې ډېر
څه ثابتوي۔ استدلال ساده و: موږ يوازې دا حقيقت وکاراوه چې تابعې
$f_0,f_1,\dots$ په داسې ډول په لړ کښې راوستل کېدے شي چې د~$x$ او~$y$
د تابع په توګه $f_x(y)$ محاسبه کېدونکې وي۔ نو دا استدلال د تابعو پر
\emph{هر هغه ټولګي عملي کېږي چې په دې ډول په لړ کښې راوستل کېدے شي}۔
له دې سره له يوې ستونزې سره مخ کېږو: غواړو محاسبه کېدونکې تابعې په
څرګند ډول بيان کړو؛ خو د محاسبه کېدونکو تابعو د يوې ټولګې هېڅ څرګند
بيان بشپړ کېدے نۀ شي!

د وتلو لار دا ده چې \emph{جزوي} تابعو ته هم ځای ورکړو۔ وبه وينو چې
جزوي محاسبه کېدونکې تابعې \emph{په رښتيا} په لړ کښې راوستل کېدے شي۔
\iftag{TMs}{په حقيقت کښې نږدې له وړاندې پوهېږو چې داسې ده، ځکه ټيورينګ
  ماشينونه په منظم ډول په لړ کښې راوستل کېدے شي۔}{} وروسته به خپل قطري
استدلال ته راوګرځو او وبه څېړو چې ولې د جزوي تابعو په ګډون هغه استدلال
مخ ته نۀ ځي۔

اوس پوښتنه دا ده: بنسټيزو بازګشتي تابعو ته څه ورزيات کړو چې ټولې جزوي
بازګشتي تابعې تر لاسه کړو؟ دوه کارونه پکار دي:
\begin{enumerate}
\item د بنسټيزو بازګشتي تابعو خپل تعريف داسې بدل کړو چې جزوي تابعو ته
  هم ځای ورکړي۔
\item تعريف ته يو څه \emph{ورزيات} کړو، څو ځينې نوې جزوي تابعې هم پکې
  شاملې شي۔
\end{enumerate}

لومړے کار اسان دے۔ د پخوا په شان له صفر، جانشين او پروجکشنونو څخه پيل
کوو، او ټولګے د ترکيب او بنسټيز بازګښت لاندې بندوو۔ يوازينے توپير دا
دے چې د ترکيب او بنسټيز بازګښت تعريفونه بايد د دې امکان دپاره بدل کړو
چې په تعريف کښې ځينې ترمونه تعريف شوي نۀ وي۔ که $f$ او~$g$ جزوي تابعې
وي، $f(x) \fdefined$ به په دې مانا ليکو چې $f$ پر~$x$ تعريف شوې ده،
يعنې $x$ د~$f$ په تعريف ساحه کښې دے؛ او $f(x) \fundefined$ به د دې
مخالفې مانا، يعنې دا چې $f$ پر~$x$ تعريف شوې نۀ ده، څرګندوي۔
$f(x) \simeq g(x)$ به په دې مانا کاروو چې يا $f(x)$ او~$g(x)$ دواړه
تعريف شوي نۀ دي، يا دواړه تعريف شوي او سره برابر دي۔ دا نښې به د لا
پېچلو ترمونو دپاره هم کاروو۔ دا دود به منو چې که $h$ او $g_0$،
\dots،~$g_k$ ټولې جزوي تابعې وي، نو
\[
h(g_0(\vec x),\dots,g_k(\vec x))
\]
هله او يوازې هله تعريف شوے دے چې هر~$g_i$ پر~$\vec x$ تعريف شوے وي،
او $h$ پر $g_0(\vec x)$، \dots،~$g_k(\vec x)$ تعريف شوې وي۔ په دې
پوهه د جزوي تابعو د ترکيب او بنسټيز بازګښت تعريفونه هماغه پورته
تعريفونه دي، خو ``$=$'' بايد پر ``$\simeq$'' بدل کړو۔

بنسټيزو بازګشتي تابعو ته چې څه ورزياتوو څو جزوي تابعې تر لاسه کړو،
هغه د \emph{بې‌حده پلټنې عامل} دے۔ که $f(x,\vec z)$ پر طبيعي عددونو
هره جزوي تابع وي، $\mu x \; f(x,\vec z)$ داسې تعريف کړئ:
\begin{quote}
  تر ټولو وړوکے $x$ چې $f(0,\vec z), f(1,\vec z), \dots, f(x,\vec
  z)$ ټول تعريف شوي وي، او $f(x,\vec z) = 0$ وي، که داسې~$x$ موجود وي؛
\end{quote}
او که داسې نۀ وي، $\mu x \; f(x,\vec z)$ تعريف شوے نۀ دے۔ دا
$\mu x \; f(x,\vec z)$ په يکتا ډول تعريفوي۔

\begin{explain}
پام وکړئ چې زموږ تعريف \iftag{TMs}{ټيورينګ ماشينونو، يا}{} الګوريتمونو
يا د محاسبې کوم ځانګړي مدل ته هېڅ اشاره نۀ کوي۔ خو د ترکيب او بنسټيز
بازګښت په شان، د بې‌حده پلټنې تر شا هم عملياتي، محاسبوي شهود شته۔ د
يوې جزوي تابع د محاسبه کېدو په بحث کښې هغه آرګومېنټونه، چې تابع پرې
تعريف شوې نۀ وي، له هغو ننوتونو سره سمون خوري چې محاسبه پرې نۀ درېږي۔
د~$\mu x \; f(x,\vec z)$ د محاسبې کړنلار داسې ده:
$f(0,\vec z), f(1,\vec z), f(2,\vec z)$ او ورپسې قيمتونه محاسبه کړئ، تر څو د
صفر قيمت راوګرځي۔ خو که له منځنيو محاسبو څخه کومه يوه ونۀ درېږي، د
$\mu x \; f(x,\vec z)$ محاسبه هم نۀ درېږي۔
\end{explain}

که $R(x,\vec z)$ هره اړيکه وي، $\mu x \; R(x,\vec z)$ داسې تعريفېږي:
$\mu x \; (1 \tsub \Char{R}(x,\vec z))$۔ په بله وينا، $\mu x \;
R(x,\vec z)$ د~$x$ تر ټولو وړوکے قيمت راګرځوي چې $R(x,\vec z)$ پوره
کوي۔ نو که $f(x,\vec z)$ هرځاے تعريف شوې تابع وي، $\mu x \;
f(x,\vec z)$ له $\mu x \; (f(x,\vec z) = 0)$ سره يو شان دے۔ خو پام
وکړئ چې زموږ لومړنی تعريف لا عمومي دے، ځکه دا امکان هم پرېږدي چې
$f(x,\vec z)$ په هر ځای کښې تعريف شوې نۀ وي (خو د دې پر خلاف د يوې
اړيکې ځانګړونکې تابع تل هرځاے تعريف شوې وي)۔

\begin{defn}
د \emph{جزوي بازګشتي تابعو} سټ له طبيعي عددونو څخه طبيعي عددونو ته د
بېلابېلو ځاييزو کچو د جزوي تابعو تر ټولو کوچنے سټ دے چې صفر، جانشين او
پروجکشنونه لري، او د ترکيب، بنسټيز بازګښت او بې‌حده پلټنې لاندې بند دے۔
\end{defn}

البته ځينې جزوي بازګشتي تابعې به په اتفاقي ډول هرځاے تعريف شوې وي،
يعنې د هر آرګومېنټ دپاره تعريف شوې وي۔

\begin{defn}
\ollabel{defn:recursive-fn}
د \emph{بازګشتي تابعو} سټ د هغو جزوي بازګشتي تابعو سټ دے چې هرځاے
تعريف شوې وي۔
\end{defn}

پر دې د ټينګار دپاره چې يوه بازګشتي تابع په هر ځای کښې تعريف شوې ده،
کله کله ورته ``هرځاے تعريف شوې بازګشتي'' تابع هم ويل کېږي۔

\end{document}
